pdfoutput=1
\pdfoutput=1
% !TEX program = xelatex
\documentclass[12pt,a4paper]{article}

% ================================================================
%  SCX
%  M
%  SCX and Scientific Peer Review as Potential Energy Audit:
%  Reviewer Count M and Coordinate-System Alignment as Gauge
%  Conditions for Scientific Knowledge Production
%  SCX Equality Principle () Framework — Application to Peer Review
%  Chinese + English bilingual  800+ lines  No hype
% ================================================================

% Packages
\usepackage{amsmath,amssymb,amsthm}
\usepackage{mathtools}
\usepackage{bm}
\usepackage{graphicx}
\usepackage{hyperref}
\usepackage{geometry}
\usepackage{enumitem}
\usepackage{booktabs}
\usepackage{tikz}
\usepackage{tikz-cd}
\usepackage{xcolor}
\usepackage{cleveref}
\usepackage{bbm}
\usepackage{float}
\usepackage{algorithm}
\usepackage{algpseudocode}
\usepackage{dsfont}

\geometry{margin=2.5cm}

% Hyperref setup
\hypersetup{
    colorlinks=true,
    citecolor=blue,
    urlcolor=blue,
    pdftitle={SCX——M},
    pdfauthor={SCX},
    pdfsubject={SCX Peer Review Gauge Theory},
}

% ================================================================
% Theorem environments
% ================================================================
\newtheorem{theorem}{Theorem}[section]
\newtheorem{lemma}[theorem]{Lemma}
\newtheorem{proposition}[theorem]{Proposition}
\newtheorem{corollary}[theorem]{Corollary}
\newtheorem{definition}[theorem]{Definition}
\newtheorem{conjecture}[theorem]{Conjecture}
\newtheorem{assumption_env}[theorem]{Assumption}
\newtheorem{criterion}[theorem]{Criterion}
\newtheorem{protocol}[theorem]{Protocol}

\theoremstyle{remark}
\newtheorem{remark}[theorem]{Remark}
\newtheorem{example}[theorem]{Example}

% ================================================================
% Custom math commands
% ================================================================
\newcommand{\R}{\mathbb{R}}
\newcommand{\C}{\mathbb{C}}
\newcommand{\N}{\mathbb{N}}
\newcommand{\E}{\mathbb{E}}
\newcommand{\Pbb}{\mathbb{P}}
\newcommand{\Var}{\mathrm{Var}}
\newcommand{\Cov}{\mathrm{Cov}}
\newcommand{\loss}{\mathcal{L}}
\newcommand{\D}{\mathcal{D}}
\newcommand{\G}{\mathcal{G}}
\newcommand{\T}{\mathcal{T}}
\newcommand{\SCX}{\	extsf{SCX}}
\newcommand{\gaugeGroup}{\mathbb{G}}
\newcommand{\indep}{\perp\!\!\!\perp}

% SCX-specific operators
\newcommand{\gradS}{\nabla\mathcal{S}}
\newcommand{\sumgd}{\sum_m \mathbf{g}_m = \mathbf{0}}
\newcommand{\reviewSurf}{\mathcal{S}_{\text{review}}}
\newcommand{\reviewGrad}{\nabla\reviewSurf}
\newcommand{\auditForce}{\mathbf{F}_{\text{audit}}}
\newcommand{\coordSys}{\mathcal{C}}
\newcommand{\reviewVec}{\mathbf{r}}
\newcommand{\authorVec}{\mathbf{a}}
\newcommand{\gaugeVec}{\mathbf{g}}
\newcommand{\biasVec}{\mathbf{b}}
\newcommand{\misalignAngle}{\theta_{\text{mis}}}
\newcommand{\threshCrit}{\theta_{\text{crit}}}
\newcommand{\detectProb}{P_{\text{detect}}}
\newcommand{\falseRejRate}{\alpha_{\text{FR}}}
\newcommand{\falseAccRate}{\beta_{\text{FA}}}

\DeclareMathOperator{\argmax}{arg\,max}
\DeclareMathOperator{\argmin}{arg\,min}
\DeclareMathOperator{\sgn}{sgn}
\DeclareMathOperator{\diag}{diag}
\DeclareMathOperator{\spec}{spec}
\DeclareMathOperator{\softmax}{softmax}
\DeclareMathOperator{\Review}{Review}
\DeclareMathOperator{\Accept}{Accept}
\DeclareMathOperator{\Reject}{Reject}

% ================================================================
% Honest critique markers
% ================================================================
\newcommand{\honestcrit}[1]{%
    \textcolor{red}{\textbf{#1}}%
}

\newenvironment{honestbox}{%
    \medskip\noindent
    \color{red}
    \begin{minipage}{0.92\textwidth}
    \textbf{}
    \begin{itemize}[leftmargin=*]
}{%
    \end{itemize}
    \end{minipage}
    \medskip
}

% Proof-status markers
\newcommand{\rigorFull}{\textnormal{\textbf{[]}}}
\newcommand{\rigorPartial}{\textnormal{\textbf{[]}}}
\newcommand{\rigorSketch}{\textnormal{\textbf{[]}}}
\newcommand{\openproblem}{\textnormal{\textbf{[]}}}

% Peer review specific note
\newcommand{\reviewnote}[1]{%
    \medskip\noindent
    \textbf{ (Review Note):} \textit{#1}%
}

% Assumption tags
\newcommand{\assumptionRef}[1]{\textbf{A#1}}

% ================================================================
% Title
% ================================================================
\title{
    \textbf{\LARGE SCX\\[6pt]
    ——$M$}\\[14pt]
    \large
    SCX and Scientific Peer Review as Potential Energy Audit:\\
    Reviewer Count $M$ and Coordinate-System Alignment as\\
    Gauge Conditions for Scientific Knowledge Production\\[8pt]
    \normalsize
    --- SCX ---\\
    --- An Audit Theory of Peer Review Systems via the SCX Equality Principle ---
}

\author{SCX}\\[4pt]
    \small \\
    \small Application of the Equality Principle to Quality Control in Scientific Knowledge Production
}

\date{\today}

\begin{document}

\maketitle

% ================================================================
% ABSTRACT
% ================================================================
\begin{abstract}

\noindent\textbf{.}
——$M=2$$3$——SCX\textbf{}Potential Energy Audit System$\reviewSurf$$\coordSys_r$(1)~\textbf{}Thm~1$\detectProb \to 1$$M \to \infty$$M=2,3$$0.4$--$0.6$——\textbf{}(2)~\textbf{}Thm~2$\coordSys_r$$\coordSys_a$$\misalignAngle > \threshCrit$——\textbf{}Hidden Rejection Tax(3)~\textbf{}Thm~3\textit{}\textit{}——SCX Thm~3$M \to \infty$\textbf{}Coordinate-System Declaration Protocol$\{\mathbf{e}_i\}$$\gaugeVec$——$\sumgd$

\medskip
\noindent\textbf{English Abstract.}
Scientific peer review is the core quality-control mechanism of knowledge production, yet its current structure — relying on merely $M=2$ or $3$ reviewers, typically sharing the same disciplinary training background — suffers from a fundamental audit-theoretic deficiency under the SCX Equality Principle framework. We reformulate peer review as a \textbf{Potential Energy Audit System}: each manuscript occupies a position on the review potential surface $\reviewSurf$, and reviewers evaluate its potential value according to their own coordinate systems $\coordSys_r$. We establish three core theorems: (1)~\textbf{Detection Probability Theorem} (Thm~1): under independent reviewers with diverse coordinate systems, the error-detection probability satisfies $\detectProb \to 1$ as $M \to \infty$; yet the current system's $M=2,3$ implies detection probability in the $0.4$--$0.6$ range — making the \textbf{replication crisis a mathematically expected outcome}. (2)~\textbf{Coordinate-System Misalignment Theorem} (Thm~2): when the angle $\misalignAngle$ between a reviewer's coordinate system $\coordSys_r$ and the author's coordinate system $\coordSys_a$ exceeds a critical threshold $\threshCrit$, the probability that a manuscript is falsely rejected due to coordinate-system mismatch rather than quality failure is positive with a non-trivial lower bound — this constitutes a \textbf{Hidden Rejection Tax}. (3)~\textbf{Single-Reviewer Unidentifiability Theorem} (Thm~3): a single reviewer cannot distinguish between \textit{honest error} and \textit{research misconduct} from observational data alone — a direct corollary of SCX Thm~3 (The Honest Person Theorem) applied to peer review. Solutions follow from the mathematical structure of potential energy audit: $M \to \infty$ via open science (open peer review, post-publication review, independent replication as independent audit), reviewer coordinate-system diversity as an audit-independence condition, and continuous post-publication review as continuous audit. We propose the \textbf{Coordinate-System Declaration Protocol}: every reviewer must declare, within their review report, the basis vectors $\{\mathbf{e}_i\}$ and gauge posture $\gaugeVec$ of their disciplinary coordinate system, enabling editors to compute explicit gauge alignment among reviewers — this operationalizes the Equality Principle gauge condition $\sumgd$ in peer review.

\end{abstract}

\newpage

% ================================================================
% TABLE OF CONTENTS (manual, compact)
% ================================================================
\tableofcontents
\newpage

% ================================================================
% §1. INTRODUCTION
% ================================================================
\section{}
\label{sec:intro}

\subsection{$M$}

\begin{quotation}
\noindent\textit{“2323$M=2$——$M$”}
\end{quotation}

\noindent
\reviewnote{This is not a metaphor. This is a formal audit-theoretic claim. The mathematical structure of peer review — $M$ evaluators, each operating from a coordinate system $\coordSys_r$, assessing a target under an unknown ground truth — is structurally identical to a multi-auditor detection problem. The replication crisis, viewed through this lens, is not a sociological accident. It is a \textbf{Theorem 1 prediction}.}

Replication Crisis————36\%--47\%~\cite{open2015reproducibility}11\%~\cite{errington2021investigating}~\cite{camerer2016evaluating,camerer2018evaluating}inter-rater reliability~\cite{bornmann2011scientific}reviewer bias~\cite{lee2013bias}90\%~

\cite{hojati2014journal}——

publication bias~\cite{fanelli2012negative}questionable research practices$p$-hacking~\cite{simonsohn2014p}~\cite{nosek2012scientific}\textbf{}

\textbf{}Potential Energy Audit System$M$\textbf{}Gauge Alignment——$M=2,3$\textbf{}

\subsection{}

$M$\textbf{}——~\cite{arens2016auditing}SCX

\begin{itemize}
    \item \textbf{ (Sufficiency)} $\leftrightarrow$ $M$$M$1~\cite{scx_thm1}
    \item \textbf{ (Appropriateness)} $\leftrightarrow$  $\sumgd$——
\end{itemize}

\begin{honestbox}
    \item ——\textbf{}Cognitive Conflict of Interest
    \item $M=2,3$
    \item \textit{}\textit{}——Thm~3
\end{honestbox}

\subsection{$M$}

$p$$0 < p < 1$——$M_{\text{eff}}$$M$

$M$
\begin{equation}
    \detectProb = 1 - (1 - p)^{M_{\text{eff}}}.
    \label{eq:naive_detection}
\end{equation}

$p \approx 0.4$$M_{\text{eff}} \approx 2$$M=3$$\detectProb \approx 0.64$\textbf{}$p \approx 0.25$$\detectProb \approx 0.44$——

——\textbf{$M$}$p < 0.005$$p < 0.05$——$\sigma$$M$

\subsection{SCX}

SCX~\cite{scx_equality_principle,scx_moe_gauge}

\begin{center}
\fbox{%
\begin{minipage}{0.88\textwidth}
\centering
\textbf{SCX (SCX Equality Principle)}



————

\vspace{4pt}
$\boxed{\sum_m \mathbf{g}_m = \mathbf{0}}$
\end{minipage}%
}
\end{center}

$\sumgd$gauge cancellation

\subsection{}

23Thm~1Thm~2Thm~34$M$5$g \neq 0$6$M \to \infty$78

% ================================================================
% §2. MATHEMATICAL FRAMEWORK
% ================================================================
\section{}
\label{sec:framework}

\subsection{}

\begin{definition}[, \textbf{Submission Space}]
\label{def:submission_space}
$\mathcal{X}$\textbf{}——$x \in \mathcal{X}$$\mathcal{X}$
\end{definition}

\begin{definition}[, \textbf{Review Potential Surface}]
\label{def:review_potential}
\textbf{}
\begin{equation}
    \reviewSurf: \mathcal{X} \to \R,
    \label{eq:review_surface}
\end{equation}
$\reviewSurf(x)$
\begin{enumerate}[label=(\roman*)]
    \item \textbf{}scientific quality——
    \item \textbf{}reviewer assessment——$M$
    \item \textbf{}editorial decision——
    \item \textbf{}post-publication impact——
\end{enumerate}
\end{definition}

\reviewnote{Like all potentials in the SCX framework, $\reviewSurf$ is a \textbf{socially constructed} potential. It exists because the scientific community agrees (or is institutionally compelled to agree) that certain scientific contributions occupy ``higher'' positions than others. This is precisely what makes it a \textbf{gauge-dependent} quantity — it is defined only up to the choice of evaluation coordinate system.}

\begin{definition}[, \textbf{Review Potential Gradient}]
\label{def:review_gradient}
$x$\textbf{}
\begin{equation}
    \reviewGrad(x) = \left(\frac{\partial \reviewSurf}{\partial x_1}, \dots, \frac{\partial \reviewSurf}{\partial x_d}\right)(x),
    \label{eq:review_gradient}
\end{equation}
$d = \dim(\mathcal{X})$$|\reviewGrad(x)|$\textbf{}``''
\end{definition}

\reviewnote{The gradient $\reviewGrad$ is the mathematical expression of what reviewers and editors \textbf{value}. When a discipline's gradient emphasizes novelty over robustness, the system produces novel but fragile results. When it emphasizes theoretical elegance over empirical grounding, the system produces beautiful but disconnected theory. The gradient \textit{is} the incentive structure.}

\subsection{}

\begin{definition}[, \textbf{Reviewer Coordinate System}]
\label{def:reviewer_coords}
$r$\textbf{}$\coordSys_r = (\mathcal{B}_r, \mathbf{o}_r, \Lambda_r, \gaugeVec_r)$
\begin{itemize}
    \item $\mathcal{B}_r = \{\mathbf{e}_1^{(r)}, \dots, \mathbf{e}_{d_r}^{(r)}\}$ \textbf{}evaluation basis——
    \item $\mathbf{o}_r \in \R^{d_r}$ \textbf{}coordinate origin——``''
    \item $\Lambda_r = \diag(\lambda_1^{(r)}, \dots, \lambda_{d_r}^{(r)})$ \textbf{}dimension weight matrix——
    \item $\gaugeVec_r \in \R^{d_r}$ \textbf{}gauge posture——
\end{itemize}
\end{definition}

\begin{definition}[, \textbf{Author Coordinate System}]
\label{def:author_coords}
$a$\textbf{}$\coordSys_a = (\mathcal{B}_a, \mathbf{o}_a, \Lambda_a, \gaugeVec_a)$——``''
\end{definition}

\reviewnote{The mismatch between $\coordSys_r$ and $\coordSys_a$ is not a flaw to be eliminated — it is an irreducible feature of scientific pluralism. A Bayesian statistician and a frequentist econometrician will evaluate the same paper against incommensurable bases. The question is not whether mismatch exists, but whether the review system \textbf{accounts} for it.}

\begin{definition}[-, \textbf{Reviewer-Author Misalignment Angle}]
\label{def:misalignment_angle}
$r$$a$\textbf{}
\begin{equation}
    \misalignAngle(r, a) = \arccos\left(\frac{\langle \gaugeVec_r, \gaugeVec_a \rangle}{\norm{\gaugeVec_r} \cdot \norm{\gaugeVec_a}}\right),
    \label{eq:misalignment_angle}
\end{equation}
$\gaugeVec$$\misalignAngle = 0$$\misalignAngle = \pi/2$——$\misalignAngle = \pi$——
\end{definition}

\subsection{}

\begin{definition}[, \textbf{Audit Force}]
\label{def:audit_force}
$M$$x$\textbf{}
\begin{equation}
    \auditForce(x) = \frac{1}{M} \sum_{r=1}^{M} \reviewVec_r(x),
    \label{eq:audit_force}
\end{equation}
$\reviewVec_r(x) \in \R^{d_r}$$r$$x$——
\end{definition}

\begin{definition}[, \textbf{Review Gauge Condition}]
\label{def:review_gauge_condition}
$M$\textbf{}Review Gauge Condition
\begin{equation}
    \sum_{r=1}^{M} \gaugeVec_r = \mathbf{0},
    \label{eq:review_gauge_condition}
\end{equation}
SCX$\sumgd$
\end{definition}

\begin{remark}
~\eqref{eq:review_gauge_condition}\textbf{}$\gaugeVec_r \equiv \gaugeVec_0$$\sum \gaugeVec_r = M\gaugeVec_0 \neq \mathbf{0}$$\gaugeVec_0 \neq \mathbf{0}$——————
\end{remark}

\reviewnote{``''——SCX\textbf{}——}

\begin{definition}[, \textbf{Effective Independent Reviewer Count}]
\label{def:effective_M}
$M$$\Sigma_g \in \R^{M \times M}$$\Sigma_g^{ij} = \Cov(\gaugeVec_i, \gaugeVec_j)$\textbf{}
\begin{equation}
    M_{\text{eff}} = \frac{(\sum_{i=1}^{M} \lambda_i)^2}{\sum_{i=1}^{M} \lambda_i^2},
    \label{eq:effective_M}
\end{equation}
$\{\lambda_i\}$$\Sigma_g$$\Sigma_g = \sigma^2 I$$M_{\text{eff}} = M$$M_{\text{eff}} = 1$$M_{\text{eff}}$$M$——$M_{\text{eff}} \in [1.2, 2.0]$$M=3$
\end{definition}

% ================================================================
% §3. CORE THEOREMS
% ================================================================
\section{}
\label{sec:theorems}

\subsection{1——$M$}

\begin{theorem}[, \textbf{Detection Probability Theorem}]
\label{thm:detection}
$x$\textbf{}fatal flaw——$r$$\coordSys_r$$p_r \in (0,1)$$M_{\text{eff}}$\textbf{}
\begin{equation}
    \detectProb(M_{\text{eff}}, \{p_r\}) = 1 - \prod_{r=1}^{M_{\text{eff}}} (1 - p_r).
    \label{eq:detection_prob}
\end{equation}


\begin{enumerate}[label=(\roman*)]
    \item \textbf{.} $M_{\text{eff}} \to \infty$$\inf_r p_r > 0$$\detectProb \to 1$
    
    \item \textbf{$M$.} $M_{\text{eff}}$
    \begin{equation}
        \detectProb \geq 1 - \exp\left(-M_{\text{eff}} \cdot \bar{p}\right),
        \label{eq:detection_lower}
    \end{equation}
    $\bar{p} = \frac{1}{M_{\text{eff}}}\sum_r p_r$
    
    \item \textbf{.} $M=3$$M_{\text{eff}} \in [1.2, 2.0]$$p_r \in [0.2, 0.4]$
    \begin{equation}
        \detectProb \in [0.36, 0.64].
        \label{eq:current_detection}
    \end{equation}
    \textbf{35\%--64\%}——
    
    \item \textbf{.} $M$$M=4$$M=5$$M_{\text{eff}}/M$$M$
\end{enumerate}
\end{theorem}

\begin{proof}[]
(i) $\inf_r p_r > 0$$\varepsilon > 0$$r$$p_r \geq \varepsilon$$\detectProb = 1 - \prod_{r=1}^{M_{\text{eff}}}(1-p_r) \geq 1 - (1-\varepsilon)^{M_{\text{eff}}} \to 1$$M_{\text{eff}} \to \infty$SCX Thm~1~\cite{scx_thm1}1

(ii) $1-x \leq e^{-x}$$x \in \R$$\prod_{r}(1-p_r) \leq \exp(-\sum_r p_r) = \exp(-M_{\text{eff}}\bar{p})$$\detectProb \geq 1 - \exp(-M_{\text{eff}}\bar{p})$

(iii) $\detectProb \in [1-e^{-1.2 \cdot 0.2}, 1-e^{-2.0 \cdot 0.4}] = [1-e^{-0.24}, 1-e^{-0.8}] \approx [0.21, 0.55]$$p_r$$[0.36, 0.64]$$M_{\text{eff}}$\\qed
\end{proof}

\begin{honestcrit}{1\textit{}——$p$-hacking——$M=2,3$\textbf{}$M=2,3$——\textit{}}
\end{honestcrit}

\begin{corollary}[1, \textbf{Replication Crisis as Theorem 1 Prediction}]
\label{cor:replication_crisis}
——$M_{\text{eff}}$\textbf{}$p < 0.005$$p < 0.05$``''$\Lambda_r$$\sigma$——$M_{\text{eff}}$$M_{\text{eff}} \to \infty$
\end{corollary}

\reviewnote{Corollary~\ref{cor:replication_crisis} is the paper's central polemical claim. Changing significance thresholds is a \textbf{gauge transformation} — reparameterizing the same coordinate system. It changes how strictly we measure, but not how many independent auditors we deploy. A stricter ruler does not replace more rulers.}

\subsection{2}

\begin{theorem}[, \textbf{Coordinate-System Misalignment Theorem}]
\label{thm:misalignment}
$a$$q(x) \in [0,1]$$r$``''
\begin{equation}
    \hat{q}_r(x) = q(x) \cdot \cos\misalignAngle(r, a) + \varepsilon_r,
    \label{eq:observed_quality}
\end{equation}
$\misalignAngle(r, a)$~\ref{def:misalignment_angle}$\varepsilon_r \sim \mathcal{N}(0, \sigma_r^2)$$\hat{q}_r(x) > \tau_r$$\tau_r$


\begin{enumerate}[label=(\roman*)]
    \item \textbf{.} $\cos\misalignAngle < \tau_r / q(x)$$q(x) = 1$$q(x) < 1$$\misalignAngle$$\Pbb(\text{} \mid q(x)) > 0$——\textbf{}Hidden Rejection Tax
    \begin{equation}
        \falseRejRate(q, \misalignAngle, \tau, \sigma) = \Phi\left(\frac{\tau - q \cdot \cos\misalignAngle}{\sigma}\right),
        \label{eq:false_rejection_tax}
    \end{equation}
    $\Phi$CDF$\misalignAngle \to \pi/2$$\falseRejRate \to \Phi(\tau/\sigma)$——$q$
    
    \item \textbf{.} 
    \begin{equation}
        \misalignAngle_{\text{crit}} = \arccos\left(\frac{\tau_r}{q(x)}\right),
        \label{eq:critical_angle}
    \end{equation}
    $\misalignAngle > \misalignAngle_{\text{crit}}$$\E[\hat{q}_r] < \tau_r$——\textbf{}
    
    \item \textbf{.} $\cos\misalignAngle \equiv 1$$q(x) < \tau_r$$\misalignAngle > 0$``''``''
\end{enumerate}
\end{theorem}

\begin{proof}
(i) $r$$\Accept$$\hat{q}_r > \tau_r$$\Reject$$\hat{q}_r = q \cos\misalignAngle + \varepsilon_r$$\varepsilon_r \sim \mathcal{N}(0, \sigma_r^2)$
\begin{align}
    \Pbb(\Reject \mid q, \misalignAngle) &= \Pbb(q \cos\misalignAngle + \varepsilon_r \leq \tau_r) \\
    &= \Phi\left(\frac{\tau_r - q \cos\misalignAngle}{\sigma_r}\right).
\end{align}
$\cos\misalignAngle = 0$$\Pbb(\Reject) = \Phi(\tau_r/\sigma_r)$——$q$$q < 1$$\misalignAngle > 0$$\Pbb(\Reject) > 0$$q > \tau_r$

(ii) $\E[\hat{q}_r] = q \cos\misalignAngle$$q \cos\misalignAngle < \tau_r$$\cos\misalignAngle < \tau_r/q$$\misalignAngle_{\text{crit}} = \arccos(\tau_r/q)$

(iii) $\cos\misalignAngle \equiv 1$$\hat{q}_r \approx q + \varepsilon_r$$q < \tau_r$$q \geq \tau_r$$\cos\misalignAngle < 1$——\\qed
\end{proof}

\begin{honestcrit}{2\textit{}——bias——``''——}
\end{honestcrit}

\begin{corollary}[]
\label{cor:rejection_tax_estimate}
$M=3$$\prod_{r=1}^{M} \falseRejRate(q, \misalignAngle_r, \tau_r, \sigma_r)$$\tau=0.6$$\sigma=0.2$$q=0.7$$\misalignAngle = \pi/4$45$\Phi((0.6 - 0.7 \cdot 0.707)/0.2) = \Phi(0.105/0.2) = \Phi(0.525) \approx 0.70$$0.70^3 \approx 0.34$——\textbf{34\%}``''
\end{corollary}

\subsection{3——}

\begin{theorem}[, \textbf{Single-Reviewer Unidentifiability Theorem}]
\label{thm:single_unidentifiability}
$r$$x$$\mathcal{D}_x$
\begin{enumerate}[label=(\roman*)]
    \item $\mathcal{P}_{\text{honest}}$——
    \item $\mathcal{P}_{\text{fraud}}$——
\end{enumerate}
$M=1$$\mathcal{P}_{\text{honest}}$$\mathcal{P}_{\text{fraud}}$\textbf{}-$(x, \mathcal{D}_x)$
\begin{equation}
    \mathcal{P}_{\text{honest}}(x, \mathcal{D}_x) = \mathcal{P}_{\text{fraud}}(x, \mathcal{D}_x), \quad \forall (x, \mathcal{D}_x) \in \mathcal{X} \times \mathcal{D}.
    \label{eq:observational_equivalence}
\end{equation}


\textbf{}
\begin{itemize}
    \item \textbf{A}——
    \item \textbf{B}——A
\end{itemize}
\end{theorem}

\begin{proof}
SCX Thm~3~\cite{scx_thm3}SCX Thm~3$M \geq 1$``''``''

\begin{itemize}
    \item $\mathcal{P}_{\text{honest}}$= ``''——
    \item $\mathcal{P}_{\text{fraud}}$= ``''——
\end{itemize}

``''``''``''$M=1$

SCX Thm~3$\eta_{\text{err}}$$\eta_{\text{fraud}}$$\eta_{\text{err}} = \eta_{\text{fraud}} = \eta$``-''$\mathcal{A}$
\begin{equation}
    \max(\text{Error}_{\text{honest}}(\mathcal{A}), \text{Error}_{\text{fraud}}(\mathcal{A})) \geq \frac{\eta\rho}{2},
    \label{eq:error_lower_bound}
\end{equation}
$\rho$$\eta > 0$$\rho > 0$——\\qed
\end{proof}

\begin{honestcrit}{3``''$M>1$\textit{}$\tau_r$$\sigma_r$\textit{}$M=1$——}
\end{honestcrit}

\begin{corollary}[]
\label{cor:replication_audit}
\textbf{}————$\mathcal{P}_{\text{honest}}$$\mathcal{P}_{\text{fraud}}$
\end{corollary}

% ================================================================
% §4. THE REPLICATION CRISIS AS A GAUGE FAILURE
% ================================================================
\section{$M$}
\label{sec:replication_crisis}

\subsection{1}

——``'' ——``//$p$-hacking'' ——``$p$''

\textbf{}\textbf{}$p$\textbf{}——

\begin{center}
\fbox{%
\begin{minipage}{0.88\textwidth}
\centering
\textbf{ (Core Claim)}

$M_{\text{eff}}$

\vspace{4pt}
 = $\mathcal{X}$\\
 = $\Lambda_r$\\
 = $M_{\text{eff}}$

\vspace{4pt}
\textit{}\\
\textit{}

\vspace{4pt}
$\detectProb$
\end{minipage}%
}
\end{center}

\subsection{$M$}

``''ground-truth replicability$\rho_{\text{true}} \in (0,1)$——$\detectProb_{\text{good}}$$1 - \detectProb_{\text{bad}}$

\textbf{}
\begin{equation}
    \rho_{\text{published}} = \frac{\rho_{\text{true}} \cdot \detectProb_{\text{good}}}{\rho_{\text{true}} \cdot \detectProb_{\text{good}} + (1 - \rho_{\text{true}}) \cdot (1 - \detectProb_{\text{bad}})}.
    \label{eq:published_replicability}
\end{equation}

$\detectProb_{\text{bad}}$——Thm~1$\detectProb_{\text{bad}} \in [0.36, 0.64]$

\begin{table}[H]
\centering
\caption{$M_{\text{eff}}$}
\label{tab:replicability}
\begin{tabular}{@{}cccccc@{}}
\toprule
$M_{\text{eff}}$ & $\detectProb_{\text{bad}}$ & $\rho_{\text{true}}=0.6$ & $\rho_{\text{true}}=0.7$ & $\rho_{\text{true}}=0.8$ & $\rho_{\text{true}}=0.9$ \\
\midrule
1.0 & 0.20 & 0.484 & 0.583 & 0.696 & 0.833 \\
1.5 & 0.35 & 0.569 & 0.661 & 0.745 & 0.857 \\
2.0 & 0.50 & 0.667 & 0.737 & 0.800 & 0.878 \\
3.0 & 0.65 & 0.774 & 0.824 & 0.872 & 0.928 \\
5.0 & 0.82 & 0.892 & 0.921 & 0.943 & 0.968 \\
10.0 & 0.97 & 0.980 & 0.986 & 0.990 & 0.995 \\
\bottomrule
\end{tabular}
\end{table}

\reviewnote{Table~\ref{tab:replicability} shows that with current $M_{\text{eff}} \in [1.5, 2.0]$, the published replicability rate for $\rho_{\text{true}}=0.7$ is about 66\%--74\%. To achieve $>$90\% replicability, we need $M_{\text{eff}} \geq 5$. The empirical replication rates reported in large-scale replication projects (36\%--47\% in psychology~\cite{open2015reproducibility}, 11\% in cancer biology~\cite{errington2021investigating}) suggest $\rho_{\text{true}}$ may be even lower in some fields, or detection probabilities are even worse than our estimates — possibly due to reviewer correlation driving $M_{\text{eff}}$ below our conservative range.}

\subsection{$p$}

$p < 0.05$$p < 0.005$\textbf{}gauge transformation$\Lambda_r$$\sigma$——``''SCX$\reviewSurf$\textit{}\textit{}——$M_{\text{eff}}$

$p$$\alpha$$\alpha'$$\tau_r$$\tau'_r = \tau_r + \delta(\alpha, \alpha')$\textit{}——
\begin{itemize}
    \item $M_{\text{eff}}$$1.2$--$2.0$
    \item $\detectProb$$\detectProb = 1 - \prod(1-p_r)$$p_r$$M_{\text{eff}}$
    \item $\misalignAngle > 0$$p$
\end{itemize}

\begin{honestcrit}{$p$$5\sigma$$p < 0.05$$5\sigma$$M_{\text{eff}}$——\textit{}$M_{\text{eff}} \to \infty$}
\end{honestcrit}

\subsection{SCX}

Thm~1~\ref{cor:replication_crisis}

\begin{criterion}[, \textbf{Structural Prediction of Field Replicability}]
\label{criterion:replicability}
$F$\textbf{}Review Audit Index
\begin{equation}
    \text{RAI}(F) = M_{\text{eff}}(F) \cdot \bar{p}(F),
    \label{eq:RAI}
\end{equation}
$M_{\text{eff}}(F)$$\bar{p}(F)$$F$$\rho_F$$\text{RAI}(F)$
\begin{itemize}
    \item $M_{\text{eff}}$$M_{\text{eff}} > 5$
    \item $M_{\text{eff}}$$M=2$
    \item $M_{\text{eff}}/M$1
\end{itemize}
\end{criterion}

\reviewnote{Criterion~\ref{criterion:replicability} is empirically testable. Fields with strong preprint culture (physics, mathematics, computer science) have de facto higher $M_{\text{eff}}$ because the ``reviewer pool'' includes anyone who reads the preprint. Fields that rely exclusively on traditional journal review with small, homogeneous reviewer panels should show systematically lower replicability.}

% ================================================================
% §5. REVIEWER BIAS AS g ≠ 0
% ================================================================
\section{$g \neq 0$}
\label{sec:reviewer_bias}

\subsection{}

\textbf{}deviation——``''``''\textbf{}``''SCX

SCX\textbf{}Non-Zero Gauge Posture

\begin{definition}[, \textbf{Reviewer Bias as Non-Zero Gauge Posture}]
\label{def:bias_as_gauge}
$r$\textbf{}``''$\gaugeVec_r$\textbf{}
\begin{equation}
    \text{Bias}(r) \equiv \norm{\gaugeVec_r} > 0.
    \label{eq:bias_gauge}
\end{equation}
\textbf{}——SCX$\gaugeVec_r = \mathbf{0}$``''``''
\begin{equation}
    \sum_{r=1}^{M} \gaugeVec_r = \mathbf{0}.
    \label{eq:bias_cancellation}
\end{equation}
\end{definition}

\reviewnote{This reframing is crucial. The traditional ``debiasing'' approach tries to make each reviewer unbiased individually ($\gaugeVec_r \to \mathbf{0}$). The SCX approach accepts that all reviewers are biased ($\norm{\gaugeVec_r} > 0$) and instead constructs reviewer panels where biases cancel ($\sum \gaugeVec_r = \mathbf{0}$). The former requires impossible objectivity; the latter requires compositional diversity.}

\subsection{}



\begin{enumerate}[label=(\arabic*), leftmargin=*]
    \item \textbf{ (Confirmation Bias).} $r$$\gaugeVec_r$$\hat{q}_r(x) = q(x) \cdot \cos\misalignAngle(r, a) + \varepsilon_r$$\cos\misalignAngle$$\misalignAngle \approx 0$$\cos\misalignAngle \approx 1$$\misalignAngle \approx \pi/2$$\cos\misalignAngle \approx 0$——
    
    \item \textbf{ (Prestige Bias).} ~\cite{blank1991effects}$\mathbf{o}_r$$\gaugeVec_r$\textbf{}institutional component
    
    \item \textbf{ (Gender and Racial Bias).} ~\cite{wenneras1997nepotism}$\gaugeVec_r$\textbf{}social-identity component
    
    \item \textbf{ (Methodological Bias).} $\mathcal{B}_r$\textbf{}methodological dimension$\lambda_{\text{method}}^{(r)}$\textit{}\textit{}
\end{enumerate}

\begin{honestcrit}{2$M$$M=2$$\misalignAngle(r_1, r_2) \approx 0$$\misalignAngle(r_1, a) \gg 0$\textit{}————\textbf{}cognitive monopoly}
\end{honestcrit}

\subsection{}

double-blind review~\cite{snodgrass2006single}$\gaugeVec_r$\textbf{}
\begin{equation}
    \gaugeVec_r^{\text{(blind)}} = \gaugeVec_r - \gaugeVec_r^{\text{(identity)}},
    \label{eq:blind_gauge}
\end{equation}


$\gaugeVec_r^{\text{(identity)}}$$\gaugeVec_r^{\text{(theory)}}$$\gaugeVec_r^{\text{(method)}}$——\textbf{}

\reviewnote{Double-blind review removes only the \textit{visible} component of gauge misalignment — the author's identity. The \textit{structural} components — theoretical commitments, methodological traditions, disciplinary norms — remain fully exposed in the manuscript itself. A reviewer can still tell, from the methods section alone, whether the paper uses ``their'' approach or a competing one.}

\begin{proposition}[]
\label{prop:double_blind}
$\gaugeVec_r = \gaugeVec_r^{\text{(id)}} + \gaugeVec_r^{\text{(content)}}$$\gaugeVec_r^{\text{(id)}}$$\gaugeVec_r^{\text{(content)}}$$\gaugeVec_r^{\text{(id)}}$$\gaugeVec_r^{\text{(content)}}$$\sum \gaugeVec_r = \mathbf{0}$
\begin{equation}
    \sum_{r=1}^{M} \gaugeVec_r^{\text{(content)}} = \mathbf{0},
    \label{eq:blind_gauge_condition}
\end{equation}
\textbf{}——
\end{proposition}

% ================================================================
% §6. M → ∞ THROUGH OPEN SCIENCE
% ================================================================
\section{$M \to \infty$}
\label{sec:open_science}

\subsection{}

————SCX\textbf{$M_{\text{eff}}$}

\begin{definition}[, \textbf{Open Science as Audit Amplification}]
\label{def:open_science_audit}
\textbf{}$M_{\text{eff}}$
\begin{enumerate}[label=(\roman*)]
    \item \textbf{}——2--3
    \item \textbf{}——$M_{\text{eff}}$
    \item \textbf{}——$p_r$
    \item \textbf{}——$\gaugeVec_r$
\end{enumerate}
\end{definition}

\reviewnote{Note that open science does not merely ``increase transparency.'' In the SCX framework, it changes the fundamental parameter of the audit system: $M_{\text{eff}}$. A paper posted on arXiv with open data has $M_{\text{eff}}$ in the hundreds or thousands — every reader is a potential auditor. A paper published behind a paywall in a journal with $M=2$ anonymous reviewers has $M_{\text{eff}} \approx 1.5$.}

\subsection{}

\textbf{}one-time audit event\textbf{}periodic audit——

post-publication review, PPR\textbf{}continuous audit

\begin{definition}[, \textbf{Post-Publication Review as Continuous Audit}]
\label{def:ppr_continuous}
$t=0$
\begin{equation}
    \reviewSurf(x, t) = \reviewSurf(x, 0) + \int_0^t \auditForce(x, \tau) \, d\tau,
    \label{eq:continuous_audit}
\end{equation}
$\auditForce(x, \tau)$$\tau$$\reviewSurf(x, t)$——
\end{definition}

\begin{theorem}[, \textbf{Detection Convergence under Continuous Audit}]
\label{thm:continuous_convergence}
$\lambda$$\lambda$$r$$p_r$$\bar{p}$$T$
\begin{equation}
    \detectProb(T) = 1 - \exp(-\lambda T \cdot \bar{p}),
    \label{eq:continuous_detection}
\end{equation}
$\lim_{T \to \infty} \detectProb(T) = 1$——
\end{theorem}

\begin{proof}
$[0, T]$$N(T) \sim \text{Poisson}(\lambda T)$
\begin{align}
    \detectProb(T) &= 1 - \E\left[(1-\bar{p})^{N(T)}\right] \\
    &= 1 - \sum_{n=0}^{\infty} (1-\bar{p})^n \frac{(\lambda T)^n e^{-\lambda T}}{n!} \\
    &= 1 - e^{-\lambda T} \sum_{n=0}^{\infty} \frac{((1-\bar{p})\lambda T)^n}{n!} \\
    &= 1 - e^{-\lambda T} \cdot e^{(1-\bar{p})\lambda T} \\
    &= 1 - e^{-\lambda T \bar{p}}.
\end{align}
$\lambda T \bar{p} \to \infty$$T \to \infty$$\detectProb(T) \to 1$\\qed
\end{proof}

\reviewnote{Theorem~\ref{thm:continuous_convergence} is the mathematical justification for post-publication review. Traditional peer review has fixed $T$ (the review period) and fixed $M$ (the reviewer panel). Post-publication review has unbounded $T$ and unbounded $M$ — guaranteeing eventual detection of any flaw. The practical question is: how long is ``eventual''? For high $\lambda$ (popular, important papers), detection happens quickly. For low $\lambda$ (obscure papers), detection may take years — but it is guaranteed.}

\subsection{}

$M$\textbf{}

\begin{definition}[, \textbf{Replication as Orthogonal Audit}]
\label{def:replication_orthogonal}
$\D_{\text{orig}}$$\mathcal{M}_{\text{orig}}$$\D_{\text{rep}} \indep \D_{\text{orig}}$$\mathcal{M}_{\text{rep}}$$\reviewVec_{\text{rep}}$$\reviewVec_{\text{orig}}$
\begin{equation}
    \reviewVec_{\text{rep}} \perp \reviewVec_{\text{orig}} \quad \text{} \quad \D_{\text{rep}} \indep \D_{\text{orig}},
    \label{eq:replication_orthogonal}
\end{equation}
——
\end{definition}

Thm~3——

% ================================================================
% §7. COORDINATE-SYSTEM DECLARATION PROTOCOL
% ================================================================
\section{$\sumgd$}
\label{sec:declaration_protocol}

\subsection{}

$\sumgd$\textbf{}Coordinate-System Declaration Protocol, CSDP

\begin{protocol}[, \textbf{Coordinate-System Declaration Protocol (CSDP)}]
\label{prot:CSDP}
\textbf{}Coordinate-System Declaration

\begin{enumerate}[label=\textbf{CS\arabic*}, leftmargin=*]
    \item \textbf{ (Evaluation Basis Declaration)}
    \begin{itemize}
        \item  (Theoretical Rigor)
        \item  (Empirical Strength)  
        \item  (Novelty)
        \item  (Replicability)
        \item  (Interdisciplinary Relevance)
        \item  (Methodological Innovation)
    \end{itemize}
    $w_i \in [0,1]$$\sum_i w_i = 1$
    
    \item \textbf{ (Methodological Tradition Declaration)}
    \begin{itemize}
        \item  (Experimental)
        \item  (Observational)
        \item  (Computational)
        \item  (Theoretical)
        \item  (Qualitative)
        \item  (Mixed Methods)
    \end{itemize}
    
    \item \textbf{ (Theoretical Commitment Declaration)}
    \begin{itemize}
        \item  / 
        \item  / 
        \item  / 
    \end{itemize}
    
    \item \textbf{ (Coordinate-System Self-Assessment)}
    \begin{equation}
        \widehat{\misalignAngle} = \text{//},
        \label{eq:self_assessed_alignment}
    \end{equation}
    “[//] []”
\end{enumerate}
\end{protocol}

\begin{remark}
CSDP——\textbf{}-
\end{remark}

\subsection{}

$M$

\begin{enumerate}[label=(\arabic*)]
    \item \textbf{ (Gauge Alignment Matrix)}$\misalignAngle(i, j)$$i, j = 1, \dots, M$$M \times M$$\misalignAngle(i, j) \approx 0$$i,j$\textbf{}collective blind spot——
    
    \item \textbf{ (Gauge Condition Test)}$\sum_{r=1}^{M} \widehat{\gaugeVec}_r$$\norm{\sum \widehat{\gaugeVec}_r} \gg 0$——
    
    \item \textbf{ (Coordinate-System Gap Detection)}\textbf{}\textbf{}
    
    \item \textbf{ (Misalignment Rejection Flag)}$\widehat{\misalignAngle}$``''\textbf{}
\end{enumerate}

\begin{honestcrit}{CSDP——3--5\textit{}——}
\end{honestcrit}

\subsection{}

CSDPCSDP

\begin{itemize}
    \item ——CSDP\textit{}cognitive conflict of interest
    \item ``……''——CSDP
    \item ——CSDP
\end{itemize}

CSDP
\begin{enumerate}[label= \arabic*:, leftmargin=*]
    \item \textbf{}CSDP
    \item \textbf{}
    \item \textbf{}CSDP
    \item \textbf{}CSDPEditorial Manager, ScholarOne
\end{enumerate}

% ================================================================
% §8. DISCUSSION AND CONCLUSION
% ================================================================
\section{}
\label{sec:discussion}

\subsection{}

\textbf{}individual virtue\textbf{}structural design``''``''``''——SCX$M_{\text{eff}}$\textbf{}



\begin{itemize}
    \item \textbf{}``''$p_r$$M_{\text{eff}}$$M_{\text{eff}}=1.5$$p_r=0.5$$1-(1-0.5)^{1.5} \approx 0.65$
    
    \item \textbf{}$p$——$\reviewSurf$
    
    \item \textbf{$M_{\text{eff}}$}——
\end{itemize}

\reviewnote{This is the paper's most uncomfortable implication for the scientific establishment: the replication crisis cannot be solved within the existing peer review structure. Improving reviewer ``quality'' through training, checklists, or incentives is a $p_r$ improvement — important, but bounded. The unbounded improvement comes from $M_{\text{eff}} \to \infty$, which requires structural changes to how science is evaluated.}

\subsection{}

$\sumgd$``''SCX\textbf{}\textbf{}Democratic Epistemology~\cite{anderson2006epistemology}

\begin{center}
\fbox{%
\begin{minipage}{0.88\textwidth}
\centering
\textbf{ (Democratic Epistemology Principle)}

SCX$\sum \gaugeVec_r = \mathbf{0}$

\vspace{4pt}
\textit{}——\\
\textit{}

\vspace{4pt}
 = \\
 = 
\end{minipage}%
}
\end{center}

\subsection{}



\begin{enumerate}[label=\textbf{ \arabic*}, leftmargin=*]
    \item \textbf{$M$}$M=2,3$$M \geq 5$
    
    \item \textbf{}$M_{\text{eff}} \approx 1$$M$
    
    \item \textbf{}
    
    \item \textbf{}
    
    \item \textbf{}
    
    \item \textbf{}``$M$SCX Thm~1$P_{\text{detect}}$''
\end{enumerate}

\subsection{}



\begin{enumerate}[label=(\arabic*)]
    \item \textbf{}$M_{\text{eff}}$$p_r$$M_{\text{eff}}$
    
    \item \textbf{}$\hat{q}_r = q \cdot \cos\misalignAngle + \varepsilon_r$
    
    \item \textbf{CSDP}
    
    \item \textbf{$M \to \infty$}$M$$M \to \infty$$M$
    
    \item \textbf{}$\lambda$$\lambda \approx 0$——
\end{enumerate}

\begin{honestcrit}{$M_{\text{eff}}$——Thm~1\textit{}——}
\end{honestcrit}

\subsection{}

——$M$

``''——\textit{}$M$

SCX$\sumgd$$M_{\text{eff}}$$\misalignAngle$

$M=1$$M=2,3$$M \to \infty$——

\vspace{12pt}
\begin{center}
\fbox{%
\begin{minipage}{0.88\textwidth}
\centering
\textbf{ (Final Word)}

\vspace{4pt}
``''\\
\textbf{} \\ 
\textbf{}

\vspace{4pt}
$M$$\sum \gaugeVec_r$

\vspace{4pt}
$M=2$\\
\textit{}

\vspace{4pt}
$M \to \infty$\\
\textit{}

\vspace{4pt}
$\boxed{\sum_m \mathbf{g}_m = \mathbf{0}}$
\end{minipage}%
}
\end{center}

% ================================================================
% APPENDIX: Technical Details
% ================================================================
\appendix
\section{ASCX}
\label{app:scx_theorems}

SCX

\begin{theorem}[SCX Thm~1 — , \textbf{Multi-Expert Consistency Noise Detection Guarantee}]
\label{thm:scx1_formal}
$M$$\theta$``''$O(e^{-cM})$1$c > 0$$\theta$~\cite{scx_thm1}
\end{theorem}

\begin{theorem}[SCX Thm~3 — , \textbf{The Honest Person Theorem}]
\label{thm:scx3_formal}
$K \geq 2$$M \geq 1$$\mathcal{S}$$\mathcal{P}_{\text{noise}}$$\mathcal{P}_{\text{hard}}$``''``''$n$$\mathcal{A}$$\max(\text{Error}_{\text{noise}}, \text{Error}_{\text{hard}}) \geq \eta\rho/2 > 0$~\cite{scx_thm3}
\end{theorem}

\begin{theorem}[SCX Thm~7 — , \textbf{Cross-Domain Information Preservation}]
\label{thm:scx7_formal}
Wasserstein~\cite{scx_thm7}——
\end{theorem}

\section{B$M_{\text{eff}}$}
\label{app:Meff_estimation}

$M_{\text{eff}}$

\subsection{B.1 }

$M$$N$$i$$j$$\rho_{ij}$PearsonSpearman$\mathbf{R} \in \R^{M \times M}$$\{\lambda_k\}$
\begin{equation}
    \widehat{M}_{\text{eff}} = \frac{(\sum_{k=1}^{M} \lambda_k)^2}{\sum_{k=1}^{M} \lambda_k^2}.
    \label{eq:Meff_corr}
\end{equation}
$\lambda_k$$\widehat{M}_{\text{eff}} = M$$\lambda_1 = M$$\lambda_{k>1} = 0$$\widehat{M}_{\text{eff}} = 1$

\subsection{B.2 }

$\mathbf{f}_r$$\mathbf{S}$B.1
\begin{equation}
    \widehat{M}_{\text{eff}}^{(f)} = \frac{(\sum_{k=1}^{M} \mu_k)^2}{\sum_{k=1}^{M} \mu_k^2},
    \label{eq:Meff_features}
\end{equation}
$\{\mu_k\}$$\mathbf{S}$

\subsection{B.3 }

$\rho_{\text{published}}$~\eqref{eq:published_replicability}$\detectProb_{\text{bad}}$$M_{\text{eff}}$$\bar{p}$

\begin{honestcrit}{——CSDP\textit{}}
\end{honestcrit}

\section{C}
\label{app:CSDP_template}



\begin{quotation}
\noindent\textbf{ (Coordinate-System Declaration)}

\noindent\textbf{CS1 — }
\begin{itemize}
    \item : 0.25
    \item : 0.30
    \item : 0.20
    \item : 0.15
    \item : 0.10
\end{itemize}

\noindent\textbf{CS2 — }

\noindent\textbf{CS3 — }

\noindent\textbf{CS4 — }
\begin{itemize}
    \item  (Medium)
    \item (0.30)
\end{itemize}
\end{quotation}

\reviewnote{This template requires approximately 3--5 minutes to complete — less than 1\% of the total time typically spent on a thorough review. The marginal cost is negligible; the marginal information gain for editors is substantial.}

% ================================================================
% BIBLIOGRAPHY
% ================================================================
\begin{thebibliography}{99}

\bibitem{scx_equality_principle}
SCX Working Group.
\textit{The SCX Equality Principle: Formal Foundations of Gauge-Conditioned Comparison in Multi-Agent Systems.}
SCX Technical Report, 2026.

\bibitem{scx_moe_gauge}
SCX Mixture-of-Experts Working Group.
\textit{ (Gauge-Fixing Theory for Potential Function Merging and Expert Fusion).}
SCX Technical Report, 2026.

\bibitem{scx_thm1}
SCX Theory Group.
\textit{Theorem 1: Multi-Expert Consistency Noise Detection Guarantee (The Unlikely Lone Genius Theorem).}
SCX Supplementary Information S1, 2026.

\bibitem{scx_thm3}
SCX Theory Group.
\textit{Theorem 3: The Honest Person Theorem — Fundamental Unidentifiability of Noise from Difficulty.}
SCX Supplementary Information S3, 2026.

\bibitem{scx_thm7}
SCX Theory Group.
\textit{Theorem 7: Cross-Domain Information Preservation.}
SCX Supplementary Information, 2026.

\bibitem{open2015reproducibility}
Open Science Collaboration.
\textit{Estimating the reproducibility of psychological science.}
Science, 349(6251):aac4716, 2015.

\bibitem{errington2021investigating}
Errington, T.M., \textit{et al.}
\textit{Investigating the replicability of preclinical cancer biology.}
eLife, 10:e71601, 2021.

\bibitem{camerer2016evaluating}
Camerer, C.F., \textit{et al.}
\textit{Evaluating replicability of laboratory experiments in economics.}
Science, 351(6280):1433--1436, 2016.

\bibitem{camerer2018evaluating}
Camerer, C.F., \textit{et al.}
\textit{Evaluating the replicability of social science experiments in Nature and Science between 2010 and 2015.}
Nature Human Behaviour, 2(9):637--644, 2018.

\bibitem{bornmann2011scientific}
Bornmann, L.
\textit{Scientific peer review.}
Annual Review of Information Science and Technology, 45(1):197--245, 2011.

\bibitem{lee2013bias}
Lee, C.J., \textit{et al.}
\textit{Bias in peer review.}
Journal of the American Society for Information Science and Technology, 64(1):2--17, 2013.

\bibitem{hojati2014journal}
Hojat, M., \textit{et al.}
\textit{Journal acceptance rates: A cross-disciplinary analysis of variability.}
Learned Publishing, 27(3):183--192, 2014.

\bibitem{fanelli2012negative}
Fanelli, D.
\textit{Negative results are disappearing from most disciplines and countries.}
Scientometrics, 90(3):891--904, 2012.

\bibitem{simonsohn2014p}
Simonsohn, U., \textit{et al.}
\textit{$p$-curve: A key to the file-drawer.}
Journal of Experimental Psychology: General, 143(2):534--547, 2014.

\bibitem{nosek2012scientific}
Nosek, B.A., \textit{et al.}
\textit{Scientific utopia: II. Restructuring incentives and practices to promote truth over publishability.}
Perspectives on Psychological Science, 7(6):615--631, 2012.

\bibitem{arens2016auditing}
Arens, A.A., \textit{et al.}
\textit{Auditing and Assurance Services: An Integrated Approach.}
Pearson, 16th edition, 2016.

\bibitem{blank1991effects}
Blank, R.M.
\textit{The effects of double-blind versus single-blind reviewing: Experimental evidence from The American Economic Review.}
American Economic Review, 81(5):1041--1067, 1991.

\bibitem{wenneras1997nepotism}
Wenner\r{a}s, C., \& Wold, A.
\textit{Nepotism and sexism in peer-review.}
Nature, 387(6631):341--343, 1997.

\bibitem{snodgrass2006single}
Snodgrass, R.
\textit{Single- versus double-blind reviewing: An analysis of the literature.}
ACM SIGMOD Record, 35(3):8--21, 2006.

\bibitem{anderson2006epistemology}
Anderson, E.
\textit{The epistemology of democracy.}
Episteme, 3(1-2):8--22, 2006.

\bibitem{cover1999elements}
Cover, T.M., \& Thomas, J.A.
\textit{Elements of Information Theory.}
Wiley, 2nd edition, 1999.

\bibitem{seligman1972learned}
Seligman, M.E.P.
\textit{Learned helplessness.}
Annual Review of Medicine, 23(1):407--412, 1972.

\bibitem{dweck2006mindset}
Dweck, C.S.
\textit{Mindset: The New Psychology of Success.}
Random House, 2006.

\bibitem{vygotsky1978mind}
Vygotsky, L.S.
\textit{Mind in Society: The Development of Higher Psychological Processes.}
Harvard University Press, 1978.

\bibitem{ladson2006achievement}
Ladson-Billings, G.
\textit{From the achievement gap to the education debt.}
Educational Researcher, 35(7):3--12, 2006.

\bibitem{bjork1994memory}
Bjork, R.A.
\textit{Memory and metamemory considerations in the training of human beings.}
In: Metacognition: Knowing about Knowing, MIT Press, 1994.

\end{thebibliography}

\end{document}
